% Route: one-player-howard.  Paste-ready LaTeX in the amsthm style of frontier.tex.
% Labels prefixed oph:.  Full file: <code_dir>/oph.tex

\subsection{Howard's rule with two actions: the dictionary, a class where it is
linear, and a family where it is not}\label{sec:oph}

\begin{proposition}[The one-player dictionary]\label{oph:dictionary}
Let $G$ be a stopping SSG with $\Vmin=\emptyset$ and $\Vmax=\{x_1,\dots,x_m\}$. For
$v\in V$ let $p^{v}$ be the law of the first visit to $\Vmax\cup\{t_0,t_1\}$ of the
token started at $v$, a datum of the graph, since only average vertices lie strictly
between controlled vertices. Then the transient Markov decision process
$M(G)$ with state set $\Vmax$, two actions per state, transition row
$p^{x_i^{(c)}}|_{\Vmax}$ and reward $q_{i,c}:=p^{x_i^{(c)}}(t_1)$ satisfies:
\begin{enumerate}[label=(\alph*),leftmargin=2.2em]
\item every row is substochastic and $\lim_kP_\sigma^k=0$ for every policy $\sigma$;
\item $\val_\sigma$ restricted to $\Vmax$ is the value of the policy $\sigma$ in
      $M(G)$, for every $\sigma$;
\item the strictly switchable set of $\sigma$ in $G$ is the set of states at which
      the other action has strictly larger $Q$-value in $M(G)$; hence the
      all-switches rule on $G$ is exactly Howard's policy-improvement rule
      (\emph{switch every improvable state}) on $M(G)$, step for step and from every
      start.
\end{enumerate}
Conversely, let $M$ be a transient MDP with $m$ states, two actions per state,
substochastic \emph{dyadic} rows of denominator $2^{d}$ and dyadic rewards, every
row leaking. Then \Cref{lem:dyadic-row} produces a stopping SSG $G(M)$ with
$\Vmin=\emptyset$, $|\Vmax|=m$ and $O(m^{2}d)$ average vertices with
$M(G(M))=M$; so Howard's runs on $M$ are the all-switches runs of $G(M)$.
\end{proposition}

\begin{proof}
Every value vector of a positional pair is harmonic at average vertices and pinned at
the sinks, so for any $\sigma$ and any $v$, $\val_\sigma(v)=\sum_{w\in\Vmax}
p^{v}_w\val_\sigma(w)+p^{v}(t_1)$: condition on the first visit to
$\Vmax\cup\{t_0,t_1\}$, which occurs almost surely because $G$ is stopping. Applying
this at $v=x_i^{(\sigma(i))}$ and using $\val_\sigma(x_i)=\val_\sigma(x_i^{(\sigma(i))})$
gives exactly the policy-evaluation equations of $M(G)$, whose solution is unique
because $I-P_\sigma$ is invertible: a $\sigma$-trap of $G$ would be a trap
(\Cref{lem:trapchar}) and $G$ is stopping. That is (a) and (b). For (c), $x_i$ is
strictly switchable at $\sigma$ iff
$\val_\sigma(x_i^{(1-\sigma(i))})>\val_\sigma(x_i^{(\sigma(i))})$, which by the same
display is the $Q$-value comparison in $M(G)$. The converse is
\Cref{lem:dyadic-row} applied to each of the $2m$ rows; the leak makes the game
stopping, and the first-passage law of the entry of the gadget for $(i,c)$ is the
prescribed row, so $M(G(M))=M$.
\end{proof}

\begin{remark}\label{oph:v3}
\Cref{oph:dictionary} is the formal version of the round-13 claim (v3) and settles it
in both directions; the ``uniform leak'' the claim asked for is the leak of
\Cref{lem:dyadic-row}. In particular the one-player half of the pivot of
\Cref{sec:refutations} \emph{is} Melekopoglou and Condon's question, with $|\Vmax|$
playing the role of the number of states; the number $N$ of vertices of the SSG is
the encoding length $O(m^{2}d)$ of the MDP, not its number of states.
\end{remark}

\begin{definition}[Sign-definite games]\label{oph:def-signdef}
Let $G$ be a stopping SSG and, for $v\in V$, let $p^{v}$ be the law of the first visit
to $C\cup\{t_0,t_1\}$ from $v$. Call $G$ \emph{sign-definite} if for every
$v\in\Vmax$ the vector $d^{v}:=\bigl(p^{v^{(1)}}-p^{v^{(0)}}\bigr)\big|_{C}$ is
sign-definite: either $d^{v}\ge0$ or $d^{v}\le0$ componentwise.
\end{definition}

\begin{theorem}[All-switches is linear on sign-definite games]\label{oph:signdef}
Let $G$ be a sign-definite stopping SSG, and let $\sigma_0,\sigma_1,\dots$ be any run
of improving switches (all-switches included). Then every $v\in\Vmax$ switches at most
twice, so the run has length at most $2|\Vmax|$; \Cref{prob:main} restricted to the
sign-definite games is decidable in polynomial time, and membership is decidable in
polynomial time.
\end{theorem}

\begin{proof}
As in \Cref{oph:dictionary}, for every Max strategy $\rho$ and every $v$,
$\val_\rho(v)=\sum_{w\in C}p^{v}_w\val_\rho(w)+p^{v}(t_1)$, this time with $C$ in
place of $\Vmax$ and $\val_\rho$ the value against a Min best response. Hence
\[
  \phi_v(t)\;:=\;\val_{\sigma_t}(v^{(1)})-\val_{\sigma_t}(v^{(0)})
  \;=\;\sum_{w\in C}d^{v}_w\,\val_{\sigma_t}(w)+e_v ,
\]
with $e_v$ independent of $t$. By \Cref{lem:switch} the vector $\val_{\sigma_t}$ is
nondecreasing in $t$, so if $d^{v}\ge0$ then $\phi_v$ is nondecreasing, and if
$d^{v}\le0$ it is nonincreasing. A vertex switches at $t$ exactly when $\phi_v(t)$ is
nonzero and its sign disagrees with the current choice; after such a switch the choice
agrees with the sign of $\phi_v$, and a monotone real sequence changes sign at most
once, so at most one further switch can occur. Hence $n_v\le2$ and, since each round
switches at least one vertex, $L\le2|\Vmax|$. The run ends at an optimal strategy
(\Cref{lem:local-global}), each round costs one Min best response and one evaluation
(\Cref{thm:eval-stopfree}), and each $p^{v}$ is one rational linear solve on the
average vertices, so both the algorithm and the membership test are polynomial.
\end{proof}

\begin{remark}\label{oph:signdef-scope}
The class is genuinely two-player and contains games with cycles and with
$a=\Theta(N)$; it is closed under \Cref{def:damping}, which multiplies every $d^{v}$
by a positive scalar, so the reduction of \Cref{thm:stopping-transform} keeps the
hypothesis. It is also easy: the content of \Cref{oph:signdef} is the contrapositive,
which is a design constraint on every attempt at a long run --- \emph{a run longer than
$2|\Vmax|$ needs a Max vertex whose two successors have first-passage laws whose
difference changes sign on $C$}, i.e.\ a vertex that genuinely trades one controlled
vertex off against another rather than trading a bundle against a constant. The
ladder $L_n$ of \Cref{def:ladder} is not sign-definite; the chain $\mathrm{RC}(k)$ of
\Cref{rem:grid-per-vertex} is. In two independent samples of random stopping SSGs
($882$ on $6$--$10$ non-sink vertices, of which $668$ are sign-definite, and $1009$ on
$8$--$13$, of which $445$ are) the longest run inside the class was $3$ and the
maximum number of switches at a single vertex was exactly $2$: the bound $n_v\le2$ is
attained.
\end{remark}

We now build a family of one-player games whose runs are longer than
$|\Vmax|$ by a constant factor. The mechanism is a \emph{freezing bias}: a copy of a
fixed block is held at its initial strategy by a bias whose size is controlled by a
single value of the previous copy, and that value crosses a threshold exactly at the
round at which the previous copy stops.

\begin{definition}[Balanced normal form]\label{oph:balanced}
A one-player normal form on states $B$, given by substochastic rows $p^{b,c}$ over
$B$ and rewards $q_{b,c}\ge0$, is \emph{balanced with mass $\rho$} if
$\sum_{b'}p^{b,c}_{b'}=\rho<1$ for every $b\in B$ and $c\in\{0,1\}$.
\end{definition}

\begin{lemma}[Balancing costs one state]\label{oph:balance}
Let $G_0$ be a nondegenerate one-player normal form on $m_0$ states with
$\rho:=\max_{b,c}\sum_{b'}p^{b,c}_{b'}<1$, all-switches run of length $h_0$ from
$\sigma^{0}$, and gaps
$g_{\min}:=\min_{\sigma,b}\bigl|A_{b,1}(y_\sigma)-A_{b,0}(y_\sigma)\bigr|>0$,
$g_{\max}$ the corresponding maximum. For every $\mu>0$ with
$\rho+\mu\max q_{b,c}<1$ and every dyadic $\varsigma$ with
$0<\varsigma<(1-\rho)\mu g_{\min}/2$ there is a \emph{balanced} nondegenerate
one-player normal form $\widehat G_0(\mu,\varsigma)$ on $m_0+1$ states with mass
$\rho$, namely: scale every reward by $\mu$; add a state $\omega$ with
$p^{\omega,c}_\omega=\rho$ and rewards $0,\varsigma$; and pad the row of $(b,c)$ with
$\rho-\sum_{b'}p^{b,c}_{b'}$ on $\omega$. Its values are
$\mu\,y_\sigma+\pi_{\sigma}\varsigma/(1-\rho)$ on the old states with
$\pi_\sigma\in[0,1]$ the pad of the chosen row, $\omega$ is never strictly switchable
at its action $1$, and the all-switches run from $(\sigma^{0},1)$ is the run of $G_0$,
of length $h_0$.
\end{lemma}

\begin{proof}
Scaling every reward by $\mu$ scales every value by $\mu$ and every action value by
$\mu$, so it changes no comparison. The added state satisfies
$x_\omega=\rho x_\omega+r_c$, i.e.\ $x_\omega=r_c/(1-\rho)$, so its action $1$ is
strictly better and it never switches from there; and the pad contributes
$\pi_{b,c}x_\omega\le\varsigma/(1-\rho)$ to an action value, so it perturbs each
comparison by at most $\varsigma/(1-\rho)<\mu g_{\min}/2$, less than half the smallest
gap. Hence every comparison keeps its sign and the improvement orientation on the old
coordinates is unchanged.
\end{proof}

\begin{theorem}[Stacking: the run adds up]\label{oph:stack}
Let $\widehat G$ be a balanced nondegenerate one-player normal form on $m_0$ states
with mass $\rho$, gaps in $[g_{\min},g_{\max}]$, and an all-switches run of length
$h_0$ from $\sigma^{0}$. For $k\ge1$ let $\mathrm{ST}(k)$ be the one-player normal form
on $k\,m_0$ states built from $k$ disjoint copies $B_1,\dots,B_k$ of $\widehat G$ by
\begin{itemize}[leftmargin=1.4em]
\item scaling the rewards of copy $j$ by $\mu_j>0$;
\item giving every state $b$ of copy $j\ge2$ the extra mass $\nu$ on a constant
      gadget of value $\Theta_j$ on the action $\sigma^{0}(b)$, and the extra mass
      $\nu$ on the state $v_{j-1}$ of copy $j-1$ on the other action, where $v_{j-1}$
      is a state that copy $j-1$ switches at the last round of its run --- the
      within-copy rows being re-padded so that the within-copy mass stays $\rho$ for
      both actions.
\end{itemize}
Let $x^{-}_j<x^{+}_j$ be the values of $v_{j-1}$ in $\mathrm{ST}(j-1)$ at the last two
strategies of its run, put $\Gamma_j:=x^{+}_j-x^{-}_j>0$, and choose
\[
  0<\epsilon<\frac{(1-\rho)g_{\min}}{(1-\rho)g_{\min}+g_{\max}},\qquad
  \Theta_j:=x^{+}_j-\epsilon\Gamma_j,\qquad
  \frac{\nu\epsilon\Gamma_j}{(1-\rho)g_{\min}}<\mu_j<\frac{\nu(1-\epsilon)\Gamma_j}{g_{\max}} .
\]
Then the all-switches run of $\mathrm{ST}(k)$ from the strategy that puts every copy at
$\sigma^{0}$ has length exactly $k\,h_0$: copy $j$ is frozen at $\sigma^{0}$ for the
first $(j-1)h_0$ rounds and performs $\widehat G$'s run during rounds
$(j-1)h_0,\dots,jh_0-1$.
\end{theorem}

\begin{proof}
Copies read only themselves and the single state $v_{j-1}$ of the previous copy, so
the transition matrix is block lower triangular and the values of copies
$1,\dots,j$ do not depend on copies $>j$; in particular $\mathrm{ST}(j-1)$'s values are
the restriction of $\mathrm{ST}(k)$'s, and the recursion defining $\Theta_j,\mu_j$ is
well founded.

Fix $j$ and a Max strategy $\sigma$ of $\mathrm{ST}(k)$, and write $x$ for the value
vector, $P_\sigma$ for the within-copy transition matrix of copy $j$ and
$y_\sigma$ for the value of $\sigma|_{B_j}$ in $\widehat G$. The equations of copy $j$
read $x_b=(P_\sigma x)_b+\mu_jq_{b,\sigma(b)}+\beta_b$ with
$\beta_b=\nu\Theta_j$ if $\sigma(b)=\sigma^{0}(b)$ and $\beta_b=\nu x_{v_{j-1}}$
otherwise. With $\delta:=x_{v_{j-1}}-\Theta_j$ and
$\chi_\sigma(b):=\mathbf 1[\sigma(b)\ne\sigma^{0}(b)]$ this is
$\beta=\nu\Theta_j\mathbf 1+\nu\delta\chi_\sigma$, and since every row of $P_\sigma$
has mass exactly $\rho$ we have $(I-P_\sigma)^{-1}\mathbf 1=(1-\rho)^{-1}\mathbf 1$,
whence
\[
  x|_{B_j}\;=\;\mu_j\,y_\sigma\;+\;\frac{\nu\Theta_j}{1-\rho}\mathbf 1
  \;+\;\nu\delta\,u_\sigma,\qquad u_\sigma:=(I-P_\sigma)^{-1}\chi_\sigma\in\Bigl[0,\tfrac1{1-\rho}\Bigr]^{B_j}.
\]
Because both rows at $b$ have the same mass $\rho$, $(p^{b,1}-p^{b,0})\cdot\mathbf 1=0$
and $(p^{b,1}-p^{b,0})\cdot u\le\rho\,\mathrm{osc}(u)$; so the comparison
$C_b:=(\text{value of action }1)-(\text{value of action }0)$ at $b$ is
\[
  C_b\;=\;\mu_j\,\Delta_b(y_\sigma)\;-\;s_b\,\nu\delta\;+\;\nu\delta\,\eta_b,
  \qquad |\eta_b|\le\frac{\rho}{1-\rho},
\]
where $\Delta_b(y):=A_{b,1}(y)-A_{b,0}(y)$ is $\widehat G$'s own difference and
$s_b=+1$ if $\sigma^{0}(b)=1$, $s_b=-1$ if $\sigma^{0}(b)=0$.

\emph{Freezing.} Suppose $\sigma|_{B_j}=\sigma^{0}$ and $x_{v_{j-1}}\le x^{-}_j$. Then
$\chi_\sigma=0$, so $\eta_b=0$ \emph{exactly}, and
$\delta\le x^{-}_j-\Theta_j=-(1-\epsilon)\Gamma_j<0$. Hence
$C_b=\mu_j\Delta_b(y_{\sigma^{0}})-s_b\nu\delta$ has the sign of $-s_b\nu\delta$,
i.e.\ the sign that keeps $\sigma^{0}(b)$, as soon as
$\nu(1-\epsilon)\Gamma_j>\mu_jg_{\max}$ --- the right-hand inequality of the display in
the statement. So no state of copy $j$ is strictly switchable.

\emph{Running.} Suppose $x_{v_{j-1}}=x^{+}_j$, so $\delta=\epsilon\Gamma_j$. Then
$|C_b-\mu_j\Delta_b(y_\sigma)|\le\nu\epsilon\Gamma_j(1+\rho/(1-\rho))
=\nu\epsilon\Gamma_j/(1-\rho)<\mu_jg_{\min}\le\mu_j|\Delta_b(y_\sigma)|$ by the
left-hand inequality; so $C_b$ and $\Delta_b(y_\sigma)$ have the same sign for every
$\sigma$ and every $b$. The improvement orientation of copy $j$ is therefore
$\widehat G$'s, and its all-switches trajectory from $\sigma^{0}$ is $\widehat G$'s run.

\emph{The induction.} By induction on $j$, copy $j$ is at $\sigma^{0}$ and frozen while
$x_{v_{j-1}}\le x^{-}_j$, which by \Cref{lem:switch} (values are nondecreasing along the
run, and $v_{j-1}$ takes the value $x^{-}_j$ at round $(j-1)h_0-1$ and $x^{+}_j$ at
round $(j-1)h_0$) is exactly the first $(j-1)h_0$ rounds; from round $(j-1)h_0$ on,
$x_{v_{j-1}}=x^{+}_j$ is constant, because copies $\le j-1$ are unaffected by copies
$\ge j$ and copy $j-1$ has reached its optimum. Copy $j$ then runs for $h_0$ rounds.
The copies therefore run one after another and the total is $k\,h_0$. Finally
$\epsilon<(1-\rho)g_{\min}/((1-\rho)g_{\min}+g_{\max})$ is exactly the condition that
the interval for $\mu_j$ is nonempty.
\end{proof}

\begin{corollary}[A one-player family with slope $12/7$]\label{oph:family}
Let $G_0$ be the $m=6$ one-player normal form behind the entry
$h^{\ast}_{1}(6)\ge12$ of \Cref{rem:four-ceilings} (round-$15$ howard-cube route; its
rows have denominator $2^{12}$ and its all-switches run from $\sigma=25$ has length
$12$), and let $\widehat G:=\widehat G_0(\mu,\varsigma)$ be its balanced form on $7$
states (\Cref{oph:balance}). Then for every $k$ there is a nondegenerate one-player
stopping SSG with $7k$ Max vertices, dyadic data of $O(k)$ bits and $O(k^{3})$ average
vertices, whose all-switches run has length $12k$. Hence
\[
  h^{\ast}_{1}(m)\;\ge\;\tfrac{12}{7}\,m\qquad\text{for } m\equiv0 \pmod 7 ,
\]
and more generally $h^{\ast}_{1}(m_1+k(m_0+1))\ge h_1+k\,h_0$ whenever a one-player
normal form on $m_1$ (resp.\ $m_0$) states has a run of length $h_1$ (resp.\ $h_0$);
in particular $h^{\ast}_{1}$ is superadditive up to one state,
$h^{\ast}_{1}(m_1+m_2+1)\ge h^{\ast}_{1}(m_1)+h^{\ast}_{1}(m_2)$.
\end{corollary}

\begin{proof}
\Cref{oph:balance} and \Cref{oph:stack}; the parameters are
$\rho=4095/4096$, $\nu=2^{-20}$, $\mu_1=2^{-12}$,
$\epsilon=g_{\min}/(2^{24}g_{\max})$, $\mu_{j+1}=2^{-31}\mu_j$ and $\Theta_{j+1}$ the
dyadic truncation of $x^{+}_{j+1}-\epsilon\Gamma_{j+1}$; each $\Theta_j$ and $\mu_j$
needs $O(j)$ bits, so \Cref{lem:dyadic-row} gives $O((7k)^2\cdot O(k))=O(k^{3})$
average vertices. For the last clause take copy $1$ to be the unpadded first game (no
bias is attached to copy $1$, so it needs no balancing) and copies $2,\dots$ to be
balanced copies of the second.
\end{proof}

\begin{remark}[What was computed]\label{oph:measured}
The family of \Cref{oph:family} was built and its run verified in exact rational
arithmetic for $k=1,\dots,10$, giving
$(m,L)=(7,12),(14,24),(21,36),(28,48),(35,60),(42,72),(49,84),(56,96),(63,108),(70,120)$;
the value vector increases pointwise at every round of every run. The raw SSGs were
also built through \Cref{lem:dyadic-row} and the runs recomputed \emph{from the games}
by a second, independent evaluator: $\mathrm{ST}(1)$ is a stopping one-player SSG on
$N=734$ vertices ($7$ Max) with a run of length $12$, and $\mathrm{ST}(2)$ one on
$N=2199$ vertices ($14$ Max) with a run of length $24$. The construction is the first
family of one-player stopping SSGs whose all-switches run exceeds $|\Vmax|$ by a
constant factor: the previous record families, $\mathrm{RC}(k)$ and the ladder, both
have $L=|\Vmax|$, and \Cref{cor:gate}'s attained example has $L=|\Vmax|+2$.
\end{remark}

\begin{proposition}[$h^{\ast}_{1}(5)\ge10$]\label{oph:h5}
Freezing the fifth Max vertex of the $m=6$ one-player game of \Cref{oph:family} at its
first action and retyping it as an average vertex (\Cref{def:subcube}) gives a
nondegenerate one-player stopping SSG on $346$ vertices with $5$ Max vertices whose
all-switches run from the strategy $01010$ has length $10$, with switched sets
\[
 \{0,1,3,4\},\{0,1,2,4\},\{0,1,4\},\{0,1,3\},\{0,1\},\{0\},\{2,3,4\},\{3,4\},\{4\},\{2\}.
\]
Its improvement orientation is an acyclic Holt--Klee unique sink orientation of the
$5$-cube of bottom-antipodal height $10$. Hence $h^{\ast}_{1}(5)\ge10$ and
$h^{\ast}_{\mathrm{LP}}(5)\ge10$, improving the entries $\ge9$ of
\Cref{rem:four-ceilings}; the Holt--Klee ceiling $h^{\ast}_{\mathrm{HK}}(5)=11$ of
\Cref{prop:hkfive} is now within one.
\end{proposition}

\begin{proof}
A finite exact computation, done twice: once on the normal form obtained by
eliminating the frozen state (stochastic complementation), and once on the explicit
SSG, by an evaluator that resolves the average part of the graph in topological order
and solves the resulting $6\times6$ rational system. All $32$ starting strategies were
tried; $10$ is the maximum, and the other eleven facets of the $m=6$ record give
$5,6,6,6,6,6,6,7,8,8,6$. The orientation was checked to be a USO, acyclic, tie-free
and Holt--Klee (the max-flow test of \Cref{prop:hkfive}).
\end{proof}

\begin{proposition}[$h^{\ast}_{1}(7)\ge12$ and $h^{\ast}_{1}(8)\ge13$ by explicit
games]\label{oph:h78}
$\mathrm{ST}(1)$ of \Cref{oph:family} is a nondegenerate one-player stopping SSG with
$7$ Max vertices ($N=734$) whose all-switches run has length $12$, so
$h^{\ast}_{1}(7)\ge12$ is attained by a game with seven Max vertices rather than by
monotonicity in $m$. A one-player normal form on $8$ states with denominator $2^{44}$,
found by annealing and verified exactly (tie-free, USO, acyclic, Holt--Klee),
has bottom-antipodal height $13$: its run from $\sigma=10011000$ switches
\[
\{0,1,2,4,6,7\},\{0,1,2,3,5\},\{0,1,2,5\},\{0,1,5\},\{0,1,3\},\{0,1\},\{0\},
\{2,3,4,5,6\},\{3,5\},\{5\},\{2\},\{4\},\{6\}.
\]
So $h^{\ast}_{1}(8)\ge13$ and $h^{\ast}_{\mathrm{LP}}(8)\ge13$. With
\Cref{oph:h5} the one-player row of \Cref{rem:four-ceilings} becomes
\[
  h^{\ast}_{1}(m)\;=\;1,\;2,\;4,\;6,\;\ge10,\;\ge12,\;\ge12,\;\ge13\qquad(m=1,\dots,8),
\]
and \Cref{oph:family} adds $h^{\ast}_{1}(7k)\ge12k$ for every $k$.
\end{proposition}

\begin{proposition}[Exhaustive census at $N\le8$]\label{oph:exhaustive}
Over \emph{all} one-player stopping SSGs whose non-sink vertices number
$n\le 6$ --- all $\binom{n+2}{2}^{\,n}$ assignments of an unordered pair of
\emph{distinct} targets to each vertex, for every split of $n$ into $|\Vmax|=m$ and
$a=n-m$; a vertex whose two successors coincide may be contracted, which lowers $n$
and does not raise $m$, so this is no loss --- the longest all-switches run over all
$2^{m}$ starting strategies is exactly $m$, for every $(m,a)$ with $3\le m\le n$
(for $m\le2$ the same bound is the known $h^{\ast}_{1}(1)=1$, $h^{\ast}_{1}(2)=2$, and
for $a=0$ it is \Cref{rem:grid-per-vertex}). In particular $L\le N-2$ on all of them,
with equality only at $a=0$, where the chain $\mathrm{RC}(m)$ attains it; average
vertices never lengthen a run at these sizes, so the smallest one-player game with
$L>|\Vmax|$ --- one exists, since $h^{\ast}_{1}(3)=4$ --- has at least $7$ non-sink
vertices. The computation is exact (integer Bareiss elimination, values compared as
numerators over the common determinant) and covers $1.5\times10^{9}$ games.
\end{proposition}

\begin{remark}[Two negative measurements]\label{oph:negative}
(i) \emph{Local translation-invariant averaging networks do not beat $m$.} Consider
one-player games in which every Max action goes to a single average vertex whose two
successors are given by a fixed local rule (one of $i-2,i-1,i,i+1,i+2$, the first or
the last state, $t_0$, $t_1$) evaluated at the state index, with one of five boundary
conventions (wrap, $t_0$, $t_1$, and the two mixed ones), and every edge damped by
$1-2^{-8}$. All $32040$ such families were evaluated at $m=4,5,6,7,8$: none has a run
longer than $m$ at $m=7$ or $m=8$. This is what \Cref{prop:serialiser} predicts for
travelling waves: a wave whose switched sets are singletons is forbidden
(\Cref{cor:isolated}), and exact cyclic symmetry forbids the switched set from
advancing by a rotation, since a rotation has finite order while the value vector
increases strictly.
(ii) \emph{Weakly coupled coprime lanes reach only $\approx1.5\,m$.} For the family in
which state $i$ of a cyclic lane has actions $\alpha x_{i-1}+a_i$ and
$\beta x_{i+1}+b_i$, and for the two-lane version in which the second action of a lane
reads a freely chosen state of the other lane, annealing over the offsets
$a_i,b_i$ (and the cross-targets) gives longest runs $8,10,11,13,13$ at
$m=5,\dots,9$ for one lane and $8,11,13,13,15$ at $m=5,7,8,9,11$ for two lanes ---
above $m$, but below the $12m/7$ of \Cref{oph:family} and far below the $pq$ that the
coprime-lane heuristic hopes for.
\end{remark}

\begin{remark}[The gap]\label{oph:gap}
\Cref{oph:family} is linear. The exact statement that would make it superlinear is:
\emph{there are one-player normal forms $G_0^{(1)},G_0^{(2)},\dots$ with
$m_0(t)\to\infty$ and $h_0(t)/m_0(t)\to\infty$} --- equivalently, by
\Cref{oph:stack}, $h^{\ast}_{1}(m)=\omega(m)$. Every value of $h^{\ast}_{1}$ known
here, $1,2,4,6,\ge10,\ge12$ for $m=1,\dots,6$, agrees with $\lfloor(m+1)^{2}/4\rfloor$
at $m\le4$ and exceeds it at $m=5$; the smallest open instance is $h^{\ast}_{1}(7)$,
where $\ge12$ is what \Cref{oph:family} gives and $16$ is what a quadratic law would
predict. What \Cref{oph:stack} shows is that the \emph{scheduling} half of a
superlinear construction is solved: because $\val(\Pi)$ increases strictly along every
run for every average mixture $\Pi$ of the states, a freezing bias controlled by one
value of the previous stage releases a block exactly at a prescribed round, at a cost
of one state per block and $O(1)$ bits of precision per block. What is missing is a
block whose own run is superlinear in its own state count, or a way of \emph{re-using}
a block --- running the same states through $\Omega(m)$ different sub-runs under
successively changed biases --- which \Cref{cor:isolated} permits only if no two of the
sub-runs repeat a switched set.
\end{remark}